\chapter{Introduction}

\section{Overview}\label{sec:intro_overview}
Networked surveillance cameras are increasingly wired directly into automated analytics workflows: person and perimeter detectors such as YOLOv8~\cite{jocher2023yolov8} and ByteTrack~\cite{zhang2022bytetrack} inspect frames and emit operational alerts, dispatch systems triage incidents across neighbouring devices, and forensic logs archive events for review. Throughout this pipeline sits an unexamined trust assumption---namely, that if an endpoint's software profile appears free of critical Common Vulnerabilities and Exposures (CVEs) and responds on its configured port, the video frames arriving from that device faithfully represent the live physical scene.

COBRA-WATCH addresses the resulting cyber-physical gap through \textbf{Integrity-Gated Dual-Channel Trust Fusion} (IG-DCTF). Instead of collapsing all trust into vulnerability counts, COBRA-WATCH separates \emph{vulnerability posture} (Channel~A, $T_{\mathrm{cve}} \in [0, 100]$) from \emph{signal authenticity} (Channel~B, $s \in [0, 1]$) and couples them through a multiplicative gate $T_{\mathrm{final}} = T_{\mathrm{cve}} \cdot G(s)$. When a feed fails physical or protocol liveness checks ($s < 0.50$), the gate floor ($G = 0.30$) caps the final score at $T_{\mathrm{final}} \le 30$ (\texttt{low\_trust}) regardless of how clean Channel~A looks. All exposure analysis operates strictly passively on aggregate search metadata, and visual analysis executes strictly offline on local or public benchmark video without ever opening unauthorized remote Real-Time Streaming Protocol (RTSP) streams.

\section{Motivation}
Internet-exposed surveillance infrastructure presents both massive cyber-posture exposure and asymmetric physical deception risk. Public scan metadata indexed by search engines such as Shodan~\cite{matherly2015shodan} and Censys routinely surfaces hundreds of thousands of unmanaged embedded camera interfaces, many vulnerable to Mirai-class credential brute-forcing~\cite{antonakakis2017mirai,griffioen2020mirai} or vendor firmware flaws~\cite{das2022kalay}. In a non-attributive aggregate snapshot collected under Censys research access (Ticket RS-469; \texttt{data/censys/2026-11-26/snapshot.json}), the query \texttt{host.location.country: "India" AND host.services.hardware.type: "CAMERA"} returned \textbf{104,218} exposed Indian camera hosts across six aggregate telemetry dimensions, with \textbf{81,285} hosts advertising unencrypted or exposed RTSP service port~554.

At the same time, security-only defenses address only half the threat surface. Suppose two identical IP cameras declare the same HTTP server banner, firmware version, and clean CVE catalog ($T_{\mathrm{cve}} = 100$). An attacker who replaces one endpoint behind a proxy with a replayed video loop, degrades optical sharpness gradually, occludes part of the lens, or proxies a lower-resolution stream leaves vulnerability metadata completely unchanged. Pure CVE scorers rate both feeds equally trustworthy, forwarding tampered alarms into automated dispatch queues. Bridging cyber-posture intelligence with signal liveness checks motivates the dual-channel IG-DCTF architecture.

\section{Problem Statement and Objectives}
\subsection*{Problem Statement}
Given heterogeneous passive camera posture metadata and optional authorized video frame streams, single-channel vulnerability trust scorers cannot detect spoofed, relayed, frozen, looped, blurred, occluded, or resolution-downgraded feeds, allowing malicious or tampered feeds from nominally secure devices to trigger high-trust alerts.

\subsection*{Objectives}
The project targets four verifiable engineering objectives within a strict passive ethics and legal boundary (India's IT Act, 2000 \S43/\S66 and DPDP Act, 2023):
\begin{enumerate}
  \item \textbf{IG-DCTF Scoring and Hard Integrity Gating:} Formulate and implement a dual-channel scoring engine with a piece-wise multiplicative gate ($\theta_{\mathrm{low}} = 0.50$, $\theta_{\mathrm{high}} = 0.85$, floor $0.30$, hard cap $30$) ensuring non-compensation for compromised signal integrity while preserving tolerance for benign single-field firmware drift.
  \item \textbf{Real-Video Tier~2 Tamper Detection (\texttt{v2}):} Build and validate \texttt{FrameLivenessTracker} capable of flagging freezes, natural-motion replay loops, abrupt/slow blur, full/partial lens occlusion, camera displacement, and resolution downgrades at real-time CPU speeds on public CUHK Avenue surveillance footage.
  \item \textbf{Passive Metadata Integration and Forensic Auditability:} Collect aggregate Censys exposure counts without storing or publishing raw target IP addresses, persist API alert decisions in an idempotent SHA-256 Merkle-chained ledger (\texttt{audit\_ledger}), and permanently prohibit active outbound host probing.
  \item \textbf{Interactive Console and Experimental Transparency:} Deliver an anonymized FastAPI + Next.js~16 operator console (\texttt{/console}, \texttt{/lab}, \texttt{/research}, \texttt{/ethics}) that enforces \texttt{/24} IP masking and explicitly separates real public metrics from labelled demo fixtures (\texttt{data/public.db}, $n = 9,320$ simulated devices) and counterfactual formula checks ($n = 120$).
\end{enumerate}

\section{Organization of the Report}
The remainder of this report follows the standard structure. Chapter~2 reviews prior literature across IoT compromise measurement, camera CVE assessment, trust-fusion models, and surveillance tamper detection, isolating the cyber-physical blind spot. Chapter~3 presents the passive threat model, IG-DCTF mathematical formulation, Tier~1/Tier~2 detectors, full-stack architecture, and software/hardware environment. Chapter~4 documents implementation workflows, automated unit/contract test suites, real Censys and CUHK Avenue evaluation results alongside CPU throughput benchmarks, and synthetic counterfactual ablation checks. Chapter~5 concludes the report with honest scope boundaries and future research directions.
